\chapter{Functions}
\chaptercontents{A & Functions, images and preimages (\S3.1) \\ B & Injections, surjections and bijections (\S3.2) \\ C & Review set 3}%
{\ess{3.1: 11, 15abc, 36; 3.E: 13, 17, 19, 34, 41}\\ \ess{3.2: 5, 8, 12, 14, 20, 40, 47; 3.E: 24, 25, 48}\\ \rec{3.1: 8, 16, 19, 33, 40, 41; 3.2: 6, 16, 19, 25, 39ab, 45, 46; 3.E: 2, 35, 37, 38, 47}\\ \bonus{3.1: 21}\\ Tested in \textbf{Quiz 2} (7 Oct).}

\hsection{Functions, images and preimages}

\begin{key}[Function]
A \term{function} $f:X\to Y$ consists of a set $X$ (the \term{domain}), a set $Y$ (the \term{codomain}) and a rule that assigns to \emph{each} $a\in X$ \emph{exactly one} element $f(a)\in Y$ (the \term{value} of $f$ at $a$). We write $a\mapsto f(a)$ for the rule.\\
A proposed rule is \term{well-defined} if (i) it produces a value for every $a\in X$, (ii) the value is unique (does not depend on how $a$ is written or on choices made), and (iii) the value lies in $Y$.\\
Two functions are \term{equal} if they have the same domain, the same codomain and $f(a)=g(a)$ for all $a\in X$. So $f:\R\to\R$, $x\mapsto x^2$ and $g:\R\to[0,\infty)$, $x\mapsto x^2$ are different functions.
\end{key}

\begin{key}[Functions you will meet constantly]
identity $\id_X:X\to X$, $a\mapsto a$; constant functions; the inclusion $U\hookrightarrow X$ for $U\subseteq X$; projections $\pi_1:X\times Y\to X$, $(a,b)\mapsto a$; the \term{characteristic function} $\chi_U:X\to\{0,1\}$ of $U\subseteq X$ ($\chi_U(a)=1$ if $a\in U$, else $0$); the floor $\lfloor\cdot\rfloor:\R\to\Z$; piecewise rules (``$f(n)=n/2$ if $n$ even, $f(n)=3n+1$ if $n$ odd''); sequences $(a_n)$ as functions $\N\to\R$; binary operations such as $+:\Z\times\Z\to\Z$.
\end{key}

\begin{hexample}[well-defined or not?]
Which of the following are well-defined functions? (a) $f:\Q\to\Z$, $f(\frac ab)=a$.\ (b) $g:\R\to\R$, $g(x)=1/x$.\ (c) $h:\Z\to\N$, $h(n)=n^2$.\ (d) $k:\R\to\Z$, $k(x)=$ the integer $n$ with $n\le x<n+1$.
\solution
(a) Not well-defined: $\frac12=\frac24$ but the rule gives $1$ and $2$: the value depends on how the input is written (fails (ii)).\\
(b) Not well-defined: no value at $x=0$ (fails (i)). It is a function $\R\setminus\{0\}\to\R$.\\
(c) Well-defined: $n^2$ is a unique natural number for every integer $n$ (iii: $n^2\ge0$).\\
(d) Well-defined: for each $x$ such an $n$ exists and is unique (if $n\le x<n+1$ and $m\le x<m+1$ then $|n-m|<1$, so $n=m$). This is $\lfloor x\rfloor$.\qed
\end{hexample}

\begin{key}[Graph, composition]
\begin{itemize}
\item The \term{graph} of $f:X\to Y$ is $\mathrm{Gr}(f)=\{(a,f(a))\mid a\in X\}\subseteq X\times Y$. A subset $G\subseteq X\times Y$ is the graph of a function iff $\forall a\in X,\ \exists!\,b\in Y,\ (a,b)\in G$ (``vertical line test'').
\item For $f:X\to Y$ and $g:Y\to Z$ the \term{composite} is $g\circ f:X\to Z$, $(g\circ f)(a)=g(f(a))$. Composition is associative, $h\circ(g\circ f)=(h\circ g)\circ f$, and $f\circ\id_X=f=\id_Y\circ f$; it is \emph{not} commutative.
\end{itemize}
\end{key}

\begin{hexample}[composition is not commutative]
Let $f,g:\R\to\R$ with $f(x)=x+1$ and $g(x)=x^2$. Compute $g\circ f$ and $f\circ g$ and show $g\circ f\neq f\circ g$.
\solution
$(g\circ f)(x)=g(x+1)=(x+1)^2$ and $(f\circ g)(x)=f(x^2)=x^2+1$. At $x=1$: $(g\circ f)(1)=4\neq2=(f\circ g)(1)$. Same domain and codomain, but the values differ at one point, so the functions are different.\qed
\end{hexample}

\begin{key}[Image and preimage]
Let $f:X\to Y$, $U\subseteq X$ and $V\subseteq Y$.
\begin{itemize}
\item The \term{image} of $U$ is $f[U]=\{f(a)\mid a\in U\}\subseteq Y$.\quad Rule: $b\in f[U]\Leftrightarrow\exists a\in U,\ f(a)=b$.
\item The \term{preimage} of $V$ is $f^{-1}[V]=\{a\in X\mid f(a)\in V\}\subseteq X$.\quad Rule: $a\in f^{-1}[V]\Leftrightarrow f(a)\in V$.
\item $f[X]$ is called \emph{the image of $f$}. Preimages exist for every function: $f^{-1}[V]$ does not require an inverse function.
\end{itemize}
\end{key}

\begin{hexample}[computing images and preimages]
Let $f:\R\to\R$, $f(x)=x^2$. Find $f[[-1,2]]$, $f^{-1}[[1,4]]$ and $f^{-1}[\{-1\}]$, with justification for the first.
\solution
$f[[-1,2]]=[0,4]$. ($\subseteq$) If $x\in[-1,2]$ then $0\le x^2\le4$. ($\supseteq$) If $y\in[0,4]$ then $\sqrt y\in[0,2]\subseteq[-1,2]$ and $f(\sqrt y)=y$.\reason{witness for the $\exists$ in the image rule}\\
$f^{-1}[[1,4]]=\{x\in\R\mid 1\le x^2\le4\}=[-2,-1]\cup[1,2]$.\qquad $f^{-1}[\{-1\}]=\{x\in\R\mid x^2=-1\}=\varnothing$.\qed
\end{hexample}

\begin{result}[Images and preimages versus set operations]
Let $f:X\to Y$, $g:Y\to Z$, $U,U'\subseteq X$ and $V,V'\subseteq Y$. Then:
\begin{center}\small
\begin{tabular}{@{}l l@{}}
\toprule
(a) $U\subseteq U'\Rightarrow f[U]\subseteq f[U']$ & (b) $V\subseteq V'\Rightarrow f^{-1}[V]\subseteq f^{-1}[V']$\\
(c) $f[U\cup U']=f[U]\cup f[U']$ & (d) $f^{-1}[V\cup V']=f^{-1}[V]\cup f^{-1}[V']$\\
(e) $f[U\cap U']\subseteq f[U]\cap f[U']$ \ (equality can fail) & (f) $f^{-1}[V\cap V']=f^{-1}[V]\cap f^{-1}[V']$\\
(g) $U\subseteq f^{-1}[f[U]]$ \ (equality can fail) & (h) $f[f^{-1}[V]]\subseteq V$ \ (equality can fail)\\
(i) $f^{-1}[Y\setminus V]=X\setminus f^{-1}[V]$ & (j) $(g\circ f)[U]=g[f[U]]$, \ $(g\circ f)^{-1}[W]=f^{-1}[g^{-1}[W]]$\\
\bottomrule
\end{tabular}
\end{center}
Preimages behave perfectly; images lose information. Equality in (e) and (g) holds for all $U,U'$ iff $f$ is injective; equality in (h) for all $V$ iff $f$ is surjective (Section B).
\end{result}

\begin{hexample}[an image inclusion and why it is not an equality]\label{ex:imgcap}
Prove that $f[U\cap U']\subseteq f[U]\cap f[U']$, and give an example where the inclusion is strict.
\solution
Let $b\in f[U\cap U']$. Fix $a\in U\cap U'$ with $f(a)=b$.\reason{image rule}\ Since $a\in U$ and $f(a)=b$, $b\in f[U]$; since $a\in U'$ and $f(a)=b$, $b\in f[U']$. Hence $b\in f[U]\cap f[U']$.\qed\\
Strict: $f:\R\to\R$, $f(x)=x^2$, $U=\{-1\}$, $U'=\{1\}$. Then $f[U\cap U']=f[\varnothing]=\varnothing$ but $f[U]\cap f[U']=\{1\}$. The proof of ``$\supseteq$'' breaks because the $a\in U$ and the $a'\in U'$ with $f(a)=f(a')=b$ need not be the same element.
\end{hexample}

\begin{hexample}[preimage of a complement]
Prove that $f^{-1}[Y\setminus V]=X\setminus f^{-1}[V]$.
\solution
For any $a$: \ $a\in f^{-1}[Y\setminus V]\Leftrightarrow a\in X\wedge f(a)\in Y\setminus V\Leftrightarrow a\in X\wedge f(a)\in Y\wedge f(a)\notin V\Leftrightarrow a\in X\wedge\neg(f(a)\in V)\Leftrightarrow a\in X\wedge a\notin f^{-1}[V]\Leftrightarrow a\in X\setminus f^{-1}[V]$, using that $f(a)\in Y$ is automatic for $a\in X$.\qed
\end{hexample}

\begin{hexample}[{$U\subseteq f^{-1}[f[U]]$}]
Prove that $U\subseteq f^{-1}[f[U]]$ for every $U\subseteq X$, and find $f$, $U$ with $U\neq f^{-1}[f[U]]$.
\solution
Let $a\in U$. Then $f(a)\in f[U]$ (witness $a$ itself), so $a\in f^{-1}[f[U]]$ by the preimage rule.\qed\\
Strict: $f(x)=x^2$ on $\R$, $U=\{1\}$: $f[U]=\{1\}$ and $f^{-1}[\{1\}]=\{-1,1\}\neq U$.
\end{hexample}

\begin{warn}
\begin{itemize}
\item $f[U]$ lives in the codomain, $f^{-1}[V]$ in the domain. Start a proof about $f[U]$ with ``Let $b\in f[U]$; fix $a\in U$ with $f(a)=b$''. Start a proof about $f^{-1}[V]$ with ``Let $a\in f^{-1}[V]$; then $f(a)\in V$''.
\item Do not write $f^{-1}(b)$ for a preimage unless $f$ is a bijection. The book's $f^{-1}[V]$ is a set.
\item A function is domain + codomain + rule. Changing the codomain changes the function (and can change surjectivity).
\end{itemize}
\end{warn}

\begin{planning}[3.1 and 3.E (first part)]
\ess{3.1: 11, 15abc, 36; 3.E: 13, 17, 19, 34, 41}\quad \rec{3.1: 8, 16, 19, 33, 40, 41; 3.E: 2, 37, 38}\quad \bonus{3.1: 21}\\[3pt]
Skills: checking well-definedness; computing composites, images and preimages of concrete sets (intervals, finite sets); proving the inclusions and equalities in the table above from the two membership rules; building counterexamples with $x\mapsto x^2$ or small finite functions.
\end{planning}

\hsection{Injections, surjections and bijections}

\begin{key}[Injective, surjective, bijective]
Let $f:X\to Y$.
\begin{itemize}
\item $f$ is \term{injective} (one-to-one) if \ $\forall a,b\in X,\ \big(f(a)=f(b)\Rightarrow a=b\big)$ \ (equivalently $a\neq b\Rightarrow f(a)\neq f(b)$).\\
  Not injective: $\exists a,b\in X,\ a\neq b\wedge f(a)=f(b)$.
\item $f$ is \term{surjective} (onto) if \ $\forall b\in Y,\ \exists a\in X,\ f(a)=b$ \ (equivalently $f[X]=Y$).\\
  Not surjective: $\exists b\in Y,\ \forall a\in X,\ f(a)\neq b$.
\item $f$ is \term{bijective} if it is both injective and surjective.
\end{itemize}
\end{key}

\begin{strategy}[Templates]
\begin{itemize}
\item \textbf{Injective:} ``Let $a,b\in X$ and suppose $f(a)=f(b)$.'' Compute until ``$a=b$''. (Do not start from $a=b$.)
\item \textbf{Not injective:} two explicit different inputs with the same output.
\item \textbf{Surjective:} ``Let $b\in Y$.'' Solve $f(a)=b$ for $a$ on scratch paper; in the proof: ``Put $a=\dots$. Then $a\in X$ \reason{check!} and $f(a)=\dots=b$.''
\item \textbf{Not surjective:} one explicit $b\in Y$ and a proof that $f(a)\neq b$ for all $a$ (often a parity or sign argument).
\item \textbf{Bijective:} both of the above, \emph{or} exhibit an inverse (see below) and verify both compositions.
\end{itemize}
\end{strategy}

\begin{hexample}[a bijection, directly]
Prove that $f:\R\to\R$, $f(x)=3x-2$ is bijective.
\solution
\emph{Injective.} Let $x,y\in\R$ with $f(x)=f(y)$. Then $3x-2=3y-2$, so $3x=3y$, so $x=y$.\\
\emph{Surjective.} Let $y\in\R$. Put $x=\frac{y+2}3\in\R$. Then $f(x)=3\cdot\frac{y+2}3-2=y$.\\
Hence $f$ is bijective.\qed
\end{hexample}

\begin{hexample}[neither injective nor surjective]
Let $g:\Z\to\Z$, $g(n)=n^2+n$. Show that $g$ is neither injective nor surjective.
\solution
\emph{Not injective:} $g(0)=0=g(-1)$ and $0\neq-1$.\\
\emph{Not surjective:} we show $1\notin g[\Z]$. Let $n\in\Z$. Then $g(n)=n(n+1)$ is even by Example~\ref{ex:kk1}, while $1$ is odd, so $g(n)\neq1$.\qed
\end{hexample}

\begin{hexample}[a bijection $\Z\to\N$ with cases]\label{ex:zigzag}
Define $h:\Z\to\N$ by $h(n)=2n$ if $n\ge0$ and $h(n)=-2n-1$ if $n<0$. Prove that $h$ is a bijection.
\solution
First, $h$ is well-defined: for $n\ge0$, $2n\in\N$; for $n<0$, $-2n-1\ge2-1=1\in\N$. Note $h(n)$ is even when $n\ge0$ and odd when $n<0$.\\
\emph{Injective.} Let $m,n\in\Z$ with $h(m)=h(n)$. Since $h(m)=h(n)$ have the same parity, $m,n$ are both $\ge0$ or both $<0$. If both $\ge0$: $2m=2n$, so $m=n$. If both $<0$: $-2m-1=-2n-1$, so $m=n$.\\
\emph{Surjective.} Let $k\in\N$. If $k$ is even, $k=2j$ with $j\in\N$, and $h(j)=2j=k$. If $k$ is odd, $k=2j+1$ with $j\in\N$; put $n=-(j+1)<0$, then $h(n)=2(j+1)-1=k$.\qed
\end{hexample}

\begin{key}[Inverses]
Let $f:X\to Y$. A function $g:Y\to X$ is
\begin{itemize}
\item a \term{left inverse} of $f$ if $g\circ f=\id_X$ \ (\ie $g(f(a))=a$ for all $a\in X$);
\item a \term{right inverse} of $f$ if $f\circ g=\id_Y$ \ (\ie $f(g(b))=b$ for all $b\in Y$);
\item an \term{inverse} of $f$ if it is both. Then $f$ has exactly one inverse, written $f^{-1}$, and $(f^{-1})^{-1}=f$.
\end{itemize}
\end{key}

\begin{result}[Inverses and the three properties]
Let $f:X\to Y$ with $X$ inhabited.
\begin{enumerate}
\item $f$ is injective $\Leftrightarrow$ $f$ has a left inverse.
\item $f$ is surjective $\Leftrightarrow$ $f$ has a right inverse.
\item $f$ is bijective $\Leftrightarrow$ $f$ has an inverse. Then $f^{-1}$ is also a bijection and $(g\circ f)^{-1}=f^{-1}\circ g^{-1}$ for bijections $f,g$.
\end{enumerate}
\end{result}

\begin{result}[Composition]
Let $f:X\to Y$ and $g:Y\to Z$.
\begin{enumerate}
\item If $f$ and $g$ are injective (surjective, bijective), then so is $g\circ f$.
\item If $g\circ f$ is injective, then $f$ is injective. \quad If $g\circ f$ is surjective, then $g$ is surjective.
\end{enumerate}
\end{result}

\begin{hexample}[injectivity of a composite]
Let $f:X\to Y$ and $g:Y\to Z$. Prove that if $g\circ f$ is injective then $f$ is injective. Show by example that $g$ need not be injective.
\solution
Suppose $g\circ f$ is injective. Let $a,b\in X$ with $f(a)=f(b)$. Applying $g$ gives $g(f(a))=g(f(b))$, \ie $(g\circ f)(a)=(g\circ f)(b)$, so $a=b$ by injectivity of $g\circ f$. Hence $f$ is injective.\qed\\
Example: $X=\{0\}$, $Y=\{0,1\}$, $Z=\{0\}$, $f(0)=0$, $g(0)=g(1)=0$. Then $g\circ f$ is injective (one-element domain) but $g$ is not.
\end{hexample}

\begin{hexample}[left inverse $\Rightarrow$ injective]
Let $f:X\to Y$ and suppose $g:Y\to X$ satisfies $g\circ f=\id_X$. Prove that $f$ is injective.
\solution
Let $a,b\in X$ with $f(a)=f(b)$. Then $a=\id_X(a)=g(f(a))=g(f(b))=\id_X(b)=b$.\qed\\
\emph{Converse idea:} if $f$ is injective and $x_0\in X$, define $g(b)=$ the unique $a$ with $f(a)=b$ when $b\in f[X]$, and $g(b)=x_0$ otherwise; then $g\circ f=\id_X$.
\end{hexample}

\begin{hexample}[bijective via an explicit inverse]
Let $f:\R\setminus\{2\}\to\R\setminus\{1\}$, $f(x)=\dfrac{x}{x-2}$. Prove that $f$ is a bijection and find $f^{-1}$.
\solution
\emph{Scratch:} solve $y=\frac x{x-2}$: $y(x-2)=x\Rightarrow x(y-1)=2y\Rightarrow x=\frac{2y}{y-1}$.\\
Define $g:\R\setminus\{1\}\to\R\setminus\{2\}$ by $g(y)=\frac{2y}{y-1}$. It is well-defined: $y\neq1$, and $g(y)=2$ would mean $2y=2y-2$, impossible, so $g(y)\neq2$.\\
For $x\neq2$: $g(f(x))=\dfrac{2\cdot\frac x{x-2}}{\frac x{x-2}-1}=\dfrac{2x}{x-(x-2)}=x$. \quad For $y\neq1$: $f(g(y))=\dfrac{\frac{2y}{y-1}}{\frac{2y}{y-1}-2}=\dfrac{2y}{2y-2(y-1)}=y$.\\
So $g\circ f=\id$ and $f\circ g=\id$; hence $f$ is a bijection with $f^{-1}=g$.\qed
\end{hexample}

\begin{hexample}[injective $\Rightarrow$ images preserve intersections]
Let $f:X\to Y$ be injective. Prove that $f[U\cap U']=f[U]\cap f[U']$ for all $U,U'\subseteq X$.
\solution
``$\subseteq$'' is Example~\ref{ex:imgcap}. ($\supseteq$) Let $b\in f[U]\cap f[U']$. Fix $a\in U$ with $f(a)=b$ and $a'\in U'$ with $f(a')=b$. Then $f(a)=f(a')$, so $a=a'$ by injectivity. Hence $a\in U\cap U'$ and $f(a)=b$, so $b\in f[U\cap U']$.\qed
\end{hexample}

\begin{warn}
\begin{itemize}
\item In an injectivity proof the hypothesis is $f(a)=f(b)$ and the goal is $a=b$. Writing ``$a=b\Rightarrow f(a)=f(b)$'' proves nothing (every function does that).
\item Surjectivity is about the \emph{codomain}: $x\mapsto x^2$ is surjective $\R\to[0,\infty)$ but not $\R\to\R$. Always check that your witness $a$ lies in the domain.
\item $f^{-1}$ as a function exists only for bijections; ``$f^{-1}(b)$'' in a proof must be preceded by ``since $f$ is bijective''.
\item For piecewise functions, injectivity needs the cross-cases too (an input from each piece); separate the pieces by a property of the outputs, as in Example~\ref{ex:zigzag}.
\end{itemize}
\end{warn}

\begin{planning}[3.2 and 3.E (second part)]
\ess{3.2: 5, 8, 12, 14, 20, 40, 47; 3.E: 24, 25, 48}\quad \rec{3.2: 6, 16, 19, 25, 39ab, 45, 46; 3.E: 35, 47}\\[3pt]
Skills: proving or refuting injectivity/surjectivity for concrete functions on $\N,\Z,\R$, products and power sets; finding inverses and verifying both compositions; proving the general facts about composites and left/right inverses; combining injectivity/surjectivity with images and preimages (equality cases of the Section A table).
\end{planning}

\hsection{Review set 3}
\begin{multicols}{2}\small
\begin{itemize}
\item Function = domain + codomain + rule; well-defined = every input, one output, in the codomain.
\item $b\in f[U]\Leftrightarrow\exists a\in U,\ f(a)=b$;\quad $a\in f^{-1}[V]\Leftrightarrow f(a)\in V$.
\item Preimages commute with $\cap,\cup,\setminus$; images only with $\cup$ (and $\subseteq$ for $\cap$).
\item Injective: $f(a)=f(b)\Rightarrow a=b$. Surjective: every $b$ has an $a$. Bijective: both $\Leftrightarrow$ inverse exists.
\item Left inverse $\Leftrightarrow$ injective; right inverse $\Leftrightarrow$ surjective.
\item $g\circ f$ injective $\Rightarrow f$ injective; $g\circ f$ surjective $\Rightarrow g$ surjective.
\item $(g\circ f)^{-1}=f^{-1}\circ g^{-1}$.
\item Counterexample factory: $x\mapsto x^2$, $n\mapsto 2n$, $n\mapsto n(n+1)$, tiny finite sets.
\end{itemize}
\end{multicols}
